\textbf{COLUNI 2022 -- Questão 35}\par\medskip Leia os textos a seguir.

\par

No final do século XIX, o imperialismo europeu na África e na Ásia se fundamentou em um ideal missionário de civilização, “[...] reforçando ainda mais os preconceitos raciais pelo sentimento de superioridade da ‘raça branca’”. Catherine Coquery-Vidrovitch. “O postulado da superioridade racial branca e da inferioridade negra.” In: Marc Ferro (org.). \textbf{O livro negro do colonialismo}. Rio de Janeiro, Ediouro, 2004, p. 766.

\par

Mais tarde, nos anos 1930, Adolf Hitler chegou ao poder na Alemanha com um discurso racial mais radical: “manter a pureza da raça, pois com o passar do tempo a mistura das raças levaria à degeneração” Enciclopédia do Holocausto. Disponível em: https://encyclopedia.ushmm.org/content/pt-br/article/victims-of-the-nazi-era-nazi-racial-ideology. Acesso em: 23 ago. 2021.

\par

Em nome da criação de uma super raça humana o regime nazista promoveu o extermínio de milhões de pessoas. O poder destrutivo dessa ideia atravessa os tempos. Nos anos 1990, Slobodan Milosevic perseguiu e matou milhares de pessoas de etnia albanesa na antiga Iugoslávia. Jamie Shea, então porta-voz da OTAN, afirmou que o mundo assistia a “um grande desastre humanitário, de uma grandeza não vista na Europa desde o final da Segunda Guerra Mundial (1939-45)”. Agências internacionais. “Otan fala em ‘genocídio’ em Kosovo”. Folha de São Paulo, 29 mar. 1999.

\par

Chegamos ao século XXI e, recentemente, o Alto Comissário das Nações Unidas para os Direitos Humanos, Zeid Ra'ad Al Hussein declarou “que o tratamento que Mianmar [país do sudeste asiático] reserva à minoria muçulmana rohingya se assemelha a um “exemplo de livro de limpeza étnica”. (France Presse. “ONU denuncia ‘limpeza étnica’ contra muçulmanos em Mianmar”. Portal G1. 11 set. 2017. Disponível em: https://g1.globo.com/mundo/noticia/onu-denuncia-um-classico-exemplo-de-limpeza-etnica-em-mianmar.ghtml. Acesso em: 23 ago. 2021.

\par

O imperialismo europeu, o racismo nazista e as limpezas étnicas na Iugoslávia em Mianmar fundamentam-se numa mesma ideia central. Assinale a afirmativa que expresse CORRETAMENTE essa ideia:\begin{enumerate}
\item Os humanos são criaturas livres, capazes de desenvolver diferentes culturas, por isso é necessário
        respeitar a diversidade e conduzir as nações a uma vida plural e ética.
\item Algumas etnias e culturas são superiores a outras, sendo legítimas a dominação e a violência.
\item A sobreposição de uma etnia em relação a outras é uma postura violenta que deve ser rechaçada por
        todos os países.
\item A cultura, tal como a vida, deve ser sempre respeitada e protegida por governos e instituições sociais.\end{enumerate}
